统一使用 vector 版本：系数下标为次数，所有精度参数 n 指保留 n 项（模 $x^n$），不再混用最高次数。固定模数 998244353，系数输入须归一化到 $[0,mod)$。每次运算的临时数组为局部变量，无需多测清空；卷积为 $O(n\log n)$，Newton 求逆／对数／指数／幂及除法为相应精度的 $O(n\log n)$。NTT 长度不超过 $2^{23}$；Ln 要求常数项为 1，Exp 要求常数项为 0，Inv 和除法的分母须可逆。
\begin{lstlisting}
namespace poly {
const ll mod=998244353;
using Poly=vector<ll>;
ll qp(ll a,ll b=mod-2){ll r=1;for(;b;b>>=1,a=a*a%mod)if(b&1)r=r*a%mod;return r;}
void NTT(Poly& a,bool inverse){
    int n=a.size();assert(n&&!(n&(n-1))&&n<=(1<<23));
    for(int i=1,j=0;i<n;i++){
        int bit=n>>1;for(;j&bit;bit>>=1)j^=bit;j^=bit;
        if(i<j)swap(a[i],a[j]);
    }
    for(int len=2;len<=n;len<<=1){
        ll root=qp(3,(mod-1)/len);if(inverse)root=qp(root);
        for(int i=0;i<n;i+=len){
            ll w=1;
            for(int j=0;j<len/2;j++,w=w*root%mod){
                ll u=a[i+j],v=a[i+j+len/2]*w%mod;
                a[i+j]=(u+v)%mod;a[i+j+len/2]=(u-v+mod)%mod;
            }
        }
    }
    if(inverse){ll inv=qp(n);for(ll &x:a)x=x*inv%mod;}
}
Poly Mul(Poly a,Poly b){
    if(a.empty()||b.empty())return {};
    int size=a.size()+b.size()-1;
    if(min(a.size(),b.size())<=32){
        Poly c(size);for(int i=0;i<int(a.size());i++)for(int j=0;j<int(b.size());j++)c[i+j]=(c[i+j]+a[i]*b[j])%mod;
        return c;
    }
    if(size>(1<<23))throw length_error("NTT capacity exceeded");
    int n=1;while(n<size)n<<=1;a.resize(n);b.resize(n);
    NTT(a,false);NTT(b,false);
    for(int i=0;i<n;i++)a[i]=a[i]*b[i]%mod;
    NTT(a,true);a.resize(size);return a;
}
Poly Inv(const Poly& a,int n=-1){
    if(n<0)n=a.size();if(!n)return {};
    assert(!a.empty()&&a[0]);Poly g{qp(a[0])};
    while(int(g.size())<n){
        int k=min(n,int(g.size())*2);Poly prefix(a.begin(),a.begin()+min(k,int(a.size())));
        Poly t=Mul(prefix,g);t.resize(k);
        for(ll &x:t)x=(mod-x)%mod;t[0]=(t[0]+2)%mod;
        g=Mul(g,t);g.resize(k);
    }
    return g;
}
Poly Dev(const Poly& a){
    if(a.size()<2)return {};
    Poly b(a.size()-1);for(int i=1;i<int(a.size());i++)b[i-1]=a[i]*i%mod;return b;
}
Poly Int(const Poly& a){
    assert(a.size()<size_t(mod));Poly b(a.size()+1),iv(a.size()+1);
    if(!a.empty())iv[1]=1;
    for(int i=2;i<=int(a.size());i++)iv[i]=(mod-mod/i)*iv[mod%i]%mod;
    for(int i=0;i<int(a.size());i++)b[i+1]=a[i]*iv[i+1]%mod;return b;
}
Poly Ln(const Poly& a,int n){
    if(!n)return {};assert(!a.empty()&&a[0]==1);
    Poly b=Mul(Dev(a),Inv(a,n));b.resize(n-1);b=Int(b);b.resize(n);return b;
}
Poly Exp(const Poly& a,int n){
    if(!n)return {};assert(a.empty()||a[0]==0);Poly g{1};
    while(int(g.size())<n){
        int k=min(n,int(g.size())*2);Poly log=Ln(g,k),t(k);
        for(int i=0;i<k;i++)t[i]=((i<int(a.size())?a[i]:0)-log[i]+mod)%mod;
        t[0]=(t[0]+1)%mod;g=Mul(g,t);g.resize(k);
    }
    return g;
}
Poly Pow(const Poly& a,unsigned long long exponent,int n){
    if(!n)return {};Poly result(n);
    if(!exponent){result[0]=1;return result;}
    int lead=0;while(lead<int(a.size())&&!a[lead])lead++;
    if(lead==int(a.size())||(__uint128_t)lead*exponent>=unsigned(n))return result;
    int shift=int((__uint128_t)lead*exponent),k=n-shift;
    ll c=a[lead],inv=qp(c);Poly unit(a.begin()+lead,a.end());unit.resize(k);
    for(ll &x:unit)x=x*inv%mod;
    Poly log=Ln(unit,k);for(ll &x:log)x=x*(exponent%mod)%mod;
    Poly b=Exp(log,k);ll scale=qp(c,exponent%(mod-1));
    for(int i=0;i<k;i++)result[i+shift]=b[i]*scale%mod;return result;
}
void trim(Poly& a){while(!a.empty()&&!a.back())a.pop_back();}
pair<Poly,Poly> Div(Poly a,Poly b){
    trim(a);trim(b);if(b.empty())throw invalid_argument("zero divisor");
    if(a.size()<b.size())return {{},a};
    int k=a.size()-b.size()+1;Poly ar=a,br=b;
    reverse(ar.begin(),ar.end());reverse(br.begin(),br.end());ar.resize(k);
    Poly q=Mul(ar,Inv(br,k));q.resize(k);reverse(q.begin(),q.end());
    Poly product=Mul(q,b);for(int i=0;i<int(product.size());i++)a[i]=(a[i]-product[i]+mod)%mod;
    a.resize(b.size()-1);trim(a);trim(q);return {q,a};
}
} // namespace poly
\end{lstlisting}

\paragraph{使用条件与多测} 接口相较旧指针版已变更：Mul(a,b)、Inv(a,n)、Ln(a,n)、Exp(a,n)、Pow(a,k,n)、Div(a,b)（返回商与余数）；输入不会被修改。零多项式在 trim/Div 中可表示为空 vector。$0^0$ 约定为 1；空精度 n=0 返回空。不同操作的中间卷积长度也须在 NTT 容量内，不能只按最终 n 判断。
